\chapter{Phase 1: Low-Fidelity Prototype and Baseline Failure Analysis}
\label{Chapter 2}

\section{Initial System Hypothesis and Architecture}
At the inception of this research, we followed the dominant paradigm in academic sign language recognition: dynamic word-level gesture classification. Our initial objective was to build a system capable of classifying 364 continuous sign language word glosses (such as \textit{``Hospital'', ``Doctor'', ``Car'', ``Election'', ``Telephone''}) directly from live video streams.

To construct a sufficiently comprehensive training corpus, we merged two prominent open-source video datasets:
\begin{itemize}
    \item \textbf{INCLUDE Dataset:} A benchmark Indian Sign Language dataset recorded across educational institutions in India, comprising isolated video clips of native signers.
    \item \textbf{WLASL (World Level American Sign Language) Dataset:} A large-scale American Sign Language dataset containing videos of diverse signers under varying illumination and camera distances.
\end{itemize}

\section{Feature Engineering: 856-Dimensional Spatio-Temporal Vectors}
Using MediaPipe Holistic, we extracted skeletal landmarks across the body pose, face, and both hands. Across consecutive video frames, we computed spatial coordinates, inter-joint Euclidean distances, and first-order velocity deltas, yielding an \textbf{856-dimensional feature vector per frame}.

To capture gesture dynamics over time, we implemented a 30-frame temporal sliding window buffer. Each input tensor $\mathbf{X} \in \mathbb{R}^{30 \times 856}$ was fed into a Spatial-Temporal Graph Convolutional Network (ST-GCN), where spatial graph edges represented biological skeletal bones and temporal graph edges connected identical joints across consecutive time steps.

\section{Empirical Failure Analysis: Why the Baseline Failed}
When we evaluated the trained ST-GCN model on live webcam feeds and held-out validation sets, the system encountered systemic failures that rendered it unusable for practical assistive communication. We identified four compounding bottlenecks:

\subsection{1. Cross-Regional Dialect Collision (Accuracy Collapse to 42.8\%)}
Sign languages are not universal; Indian Sign Language (ISL) and American Sign Language (ASL) possess completely distinct gestural vocabularies and syntactic conventions. For instance, the sign for \textit{``Hospital''} in ISL involves forming a two-handed cross near the upper shoulder, whereas in ASL it requires drawing an \textit{``H''} handshape downward across the opposite arm.

By merging the INCLUDE and WLASL datasets into a single 364-class taxonomy, the neural network was forced to learn contradictory spatio-temporal trajectories mapped to the exact same target labels. Consequently, the model suffered from severe catastrophic interference, achieving a dismal \textbf{42.8\% Top-1 classification accuracy} on validation test splits.

\subsection{2. Intolerable Temporal Latency ($>500\text{ ms}$ Buffer Delay)}
A dynamic sequence classifier requiring 30 video frames at 30 FPS cannot compute a prediction until the full 30-frame window is populated. This introduced a mandatory \textbf{1,000-millisecond physical gesture buffering delay}, followed by 80--120ms of neural inference time. In live conversational settings, signers felt an unnatural lag between physical hand motion and on-screen feedback, disrupting the rhythm of communication.

\subsection{3. The Closed-Dictionary Barrier}
Word-level models are strictly bounded by their training vocabulary. When a user attempted to sign an Indian personal name (e.g., \textit{``Naseer''} or \textit{``Amaan''}), a specific local town, or a specialized technical term not present in the 364-word training dictionary, the system failed completely, outputting arbitrary high-loss classifications.

\subsection{4. Single-Threaded GUI Freezing (Frame Drops to 5--12 FPS)}
Our initial desktop implementation executed OpenCV video capture, MediaPipe tracking, PyTorch neural inference, and Tkinter GUI rendering in a single sequential execution thread. The computational cost of forward-passing 30-frame temporal tensors through the ST-GCN starved the video capture loop of CPU cycles, causing the camera feed to stutter severely from 30 FPS down to \textbf{5--12 FPS}.

\begin{table}[htbp]
\centering
\caption{Summary of Baseline ST-GCN Limitations}
\label{tab:baseline_limitations}
\begin{tabular}{@{}lll@{}}
\toprule
\textbf{Dimension} & \textbf{Low-Fidelity Baseline} & \textbf{Observed Impact} \\
\midrule
Recognition Paradigm & 30-Frame Word Sequences & 1,000ms buffer lag before prediction \\
Vocabulary Reach & Closed 364-Word Dictionary & Complete failure on proper nouns and names \\
Dataset Integration & Merged ISL + ASL Videos & Severe dialect clashes (42.8\% accuracy) \\
System Architecture & Single-Threaded Python Script & Severe GUI freezing (5--12 FPS) \\
\bottomrule
\end{tabular}
\end{table}

This empirical failure forced a fundamental reassessment of our technical approach, prompting the strategic architectural pivot detailed in Chapter~3.
